\section{Methodology}

\subsection{Overview}

The proposed framework is an end-to-end pipeline for automatic leak detection
from multivariate time-series sensor data. It consists of four main stages:
(i) data preprocessing and feature engineering, (ii) sequence construction using
a sliding-window approach, (iii) classification using a hybrid CNN-BiLSTM-Attention
architecture, and (iv) threshold calibration for operational decision-making.

The same preprocessing and data splitting strategy is applied to all models
considered in this study to ensure fair comparison in the benchmark evaluation.

\subsection{Sliding Window Formulation}

Given a multivariate time series $\mathbf{X} = \{x_1, x_2, \dots, x_T\}$, where
$x_t \in \mathbb{R}^d$ represents the feature vector at time $t$, the data is
transformed into supervised learning samples using a sliding window of size $W$
and stride $S$.

Each input sequence is defined as:

\begin{equation}
\mathbf{X}_t = \{x_t, x_{t+1}, \dots, x_{t+W-1}\}
\end{equation}

with corresponding label:

\begin{equation}
y_t \in \{0,1\}
\end{equation}

where $y_t = 1$ indicates the presence of at least one leak event within the window.

In this work, a window size of $W = 20$ and stride $S = 5$ are used, allowing
the model to capture short-term temporal patterns while preserving sufficient
overlap between consecutive sequences.

\subsection{Feature Engineering}

To enhance the representation of raw sensor signals, domain-informed features
are constructed, including:

\begin{itemize}
    \item Pressure gradient ($\Delta P_t$) to capture abrupt pressure variations,
    \item Flow rate normalization (z-score) to highlight relative anomalies,
    \item Pressure-to-flow ratio to model hydraulic relationships.
\end{itemize}

These features are concatenated with the original variables to form enriched
feature vectors at each time step.

\subsection{Model Architecture}

The proposed model integrates convolutional, recurrent, and attention-based
components to capture both local and global temporal dependencies.

\subsubsection{Convolutional Layers}

The input sequences are first processed by one-dimensional convolutional layers
to extract local temporal patterns:

\begin{equation}
h^{(c)} = \text{ReLU}(\text{Conv1D}(\mathbf{X}))
\end{equation}

These layers act as feature extractors, capturing short-term variations such as
pressure drops and flow spikes.

\subsubsection{Bidirectional LSTM}

The extracted features are then passed to a Bidirectional LSTM (BiLSTM), which
captures long-range dependencies in both forward and backward directions:

\begin{equation}
h^{(l)} = \text{BiLSTM}(h^{(c)})
\end{equation}

This enables the model to incorporate contextual information from both past and
future timesteps within each sequence.

\subsubsection{Self-Attention Mechanism}

To emphasize the most informative time steps, a self-attention mechanism is
applied to the BiLSTM outputs:

\begin{equation}
\alpha_t = \frac{\exp(e_t)}{\sum_{k} \exp(e_k)}, \quad
e_t = \mathbf{v}^\top \tanh(\mathbf{W} h^{(l)}_t)
\end{equation}

The resulting context vector is computed as:

\begin{equation}
c = \sum_t \alpha_t h^{(l)}_t
\end{equation}

This mechanism allows the model to focus on temporally localized leak signatures
while preserving global sequence information.

\subsubsection{Output Layer}

The context vector is passed through a fully connected layer with sigmoid
activation to produce the final prediction:

\begin{equation}
\hat{y} = \sigma(\mathbf{W}_o c + b)
\end{equation}

where $\hat{y} \in [0,1]$ represents the probability of a leak.

\subsection{Training Strategy}

The model is trained using binary cross-entropy loss:

\begin{equation}
\mathcal{L} = - \frac{1}{N} \sum_{i=1}^{N} \left[
y_i \log(\hat{y}_i) + (1 - y_i)\log(1 - \hat{y}_i)
\right]
\end{equation}

The dataset is split into training, calibration, and test sets using temporal
separation to prevent data leakage. Early stopping is applied based on
validation AUC-PR, and learning rate scheduling is used to improve convergence.

All baseline models are trained using the same data splits and preprocessing
pipeline to ensure a fair benchmark comparison.

\subsection{Threshold Calibration}

Instead of using a fixed threshold (e.g., 0.5), the decision threshold is
optimized on a calibration set to balance precision and recall:

\begin{equation}
\hat{y}_{final} =
\begin{cases}
1 & \text{if } \hat{y} \geq \tau \\
0 & \text{otherwise}
\end{cases}
\end{equation}

This enables the model to operate under different trade-offs depending on
application requirements, such as prioritizing detection sensitivity or
minimizing false alarms.

\subsection{Implementation Details}

The model is implemented using a deep learning framework with a compact
architecture comprising approximately 23{,}000 trainable parameters.
Optimization is performed using the Adam optimizer, and training is conducted
using mini-batch gradient descent over multiple epochs.

All experiments are conducted under consistent settings to ensure reproducibility
and fair comparison across different models.